\section{Related Work}
\label{sec:related_work}

\subsection{Panoramic Semantic Segmentation and Domain Adaptation}
Panoramic semantic segmentation requires pixel-level prediction across complete surrounding fields of view~\cite{ma2021densepass,yang2021omnirange}. ERP representation facilitates processing with 2D networks but introduces stretching in high-latitude regions and separates geometrically continuous left-right boundaries; objects spanning multiple fields of view require models to simultaneously capture local details and global structures. Existing work mainly improves panoramic representation through distortion-aware convolution, projection correction, multi-view unfolding, and local-global context fusion~\cite{zhang2022bending,zheng2023dppass,zheng2023datr}. These designs focus on adapting to panoramic geometry, while this paper studies memory computation efficiency within existing panoramic frameworks.

Unsupervised domain adaptation for semantic segmentation leverages labeled source domains and unlabeled target domains to reduce distribution gaps~\cite{tsai2018adaptsegnet,zou2018cbst}. Feature or output alignment aims to learn domain-invariant representations, while self-training and consistency constraints utilize target domain predictions to construct supervision~\cite{hoyer2022daformer,hoyer2022hrda,hoyer2023mic}. Cityscapes $\rightarrow$ DensePASS involves not only appearance changes but also structural changes from perspective imaging to panoramic projection~\cite{zhang2024goodsam,zhong2025omnisam}; pseudo-labels for boundaries, distorted regions, and rare classes are thus more error-prone. This paper follows OmniSAM's self-training pipeline without redesigning complete domain alignment frameworks, instead focusing on memory token reduction and auxiliary analysis of target domain supervision reliability. Compared to approaches that only optimize feature distribution alignment, this paper focuses more on computational redundancy arising from panoramic sequence modeling in existing adaptation pipelines.

\subsection{SAM/SAM2 and Memory-based Visual Modeling}
SAM~\cite{kirillov2023segment} series foundation models provide strong general visual representations, with SAM2~\cite{ravi2025sam2} style memory mechanisms further supporting cross-frame, cross-view, and cross-patch information propagation. For panoramic images divided into local windows, historical memory can connect large-scale structures such as roads, buildings, and sky, reducing discontinuities caused by independent window predictions. OmniSAM~\cite{zhong2025omnisam} introduces this mechanism to panoramic UDA, enabling current views to leverage context from surrounding fields of view.

The cost of memory-based modeling is that attention scale grows with the number of historical tokens. In long sequences, repetitive backgrounds, adjacent regions, and weakly relevant historical representations increase training computation and may limit efficiency of ablation and parameter search. This paper keeps the SAM2-style encoder and OmniSAM segmentation pipeline unchanged, only compressing and retrieving historical representations before entering memory attention, to investigate whether cross-field context can maintain major semantic structures under lower token budgets. This problem differs from shortening original input sequences: what is reduced is historical memory, while visual features and output resolution of current views remain unchanged.

\subsection{Token-efficient Modeling and Pseudo-label Reliability}
In vision Transformers, token pruning, token merging, sparse attention, and memory compression often reduce redundant computation based on importance, similarity, or structural priors~\cite{rao2021dynamicvit,bolya2023tome,lu2023tokensharing}. Related research mainly addresses general forward inference, classification, or video long sequences; this paper addresses memory tokens in OmniSAM UDA training. MTC compresses local redundancy through spatial aggregation, while SMR retains more relevant historical tokens through importance ranking. Together they reduce the input scale to memory attention, without claiming applicability to all vision tokens or verified deployment-time acceleration.

Pseudo-labels are the core supervision source for UDA self-training, and their errors accumulate during training~\cite{zou2018cbst,hoyer2023mic}. Confidence thresholds, class balancing, uncertainty estimation, region consistency, and prototype similarity can all filter target domain predictions. Uniform thresholds easily bias toward easily classified categories, while prototype guidance relies on reliability of feature space and class centers~\cite{zheng2024sfuda360,zheng2025sfuda360pp}. This paper examines these factors through class balancing and prototype gating, but experimental gains are small, so PQR serves only as auxiliary exploration, with MTC+SMR remaining the sole method mainline. Overall, this paper connects panoramic cross-field memory with token-efficient training: the former defines context that needs retention, while the latter controls training costs paid to obtain that context.
